%%%%%%%%%%%%%%%%%%%%%%%%%%%%%%%%%%%%%%%%%%%%%%%%%%%
%% Exercices — Informatique Quantique — Chapitre 4
%% Stage LIA / CERI — Université d'Avignon
%% Stagiaire : Patrice SEBASTIANO
%% Encadrant : Serigne GUEYE
%%%%%%%%%%%%%%%%%%%%%%%%%%%%%%%%%%%%%%%%%%%%%%%%%%%

% \documentclass[light]{ceri/sty/rapport}
% \documentclass[handout]{ceri/sty/rapport}
\documentclass{ceri/sty/rapport}


%%%%%%%%%%%%%%%%%%%%%%%%%%%%%%%%%%%%%%%%%%%%%%%%%%%
%% Informations générales
%%%%%%%%%%%%%%%%%%%%%%%%%%%%%%%%%%%%%%%%%%%%%%%%%%%
\major{Master d'Informatique}
\specialization{Informatique Quantique}
\course{Stage de Recherche — LIA Avignon}

\title{Exercices — Chapitre 4}
\subtitle{Optimisation combinatoire quantique}

\author{
	Patrice SEBASTIANO
}

\advisor[Encadrant]{
	Serigne GUEYE
}


%%%%%%%%%%%%%%%%%%%%%%%%%%%%%%%%%%%%%%%%%%%%%%%%%%%
%% Packages supplémentaires
%%%%%%%%%%%%%%%%%%%%%%%%%%%%%%%%%%%%%%%%%%%%%%%%%%%
\usepackage{ifthen}
\usepackage{amsmath, amssymb}
\usepackage{mathtools}
\usepackage{braket}
\usepackage{tcolorbox}
\tcbuselibrary{breakable, skins, listings}
\usepackage{listings}
\lstset{
  language=Python,
  basicstyle=\ttfamily\small,
  keywordstyle=\bfseries\color{blue!70!black},
  stringstyle=\color{green!50!black},
  commentstyle=\itshape\color{gray},
  breaklines=true,
  frame=none,
  showstringspaces=false,
}

% Couleurs personnalisées
\definecolor{exo-bg}{RGB}{230, 240, 255}
\definecolor{exo-border}{RGB}{50, 100, 200}
\definecolor{sol-bg}{RGB}{230, 255, 235}
\definecolor{sol-border}{RGB}{40, 160, 60}

% Dessin de graphes
\usetikzlibrary{graphs, graphdrawing}
\usegdlibrary{force}

\tikzset{
  vertex/.style={circle, draw=gray!70, fill=gray!25, thick, minimum size=22pt, font=\bfseries\small},
  vset1/.style={circle, draw=blue!70, fill=blue!20, thick, minimum size=22pt, font=\bfseries\small},
  vset2/.style={circle, draw=red!70,  fill=red!20,  thick, minimum size=22pt, font=\bfseries\small},
  cut edge/.style={dashed, thick, draw=black!60},
  kept edge/.style={thick, draw=gray!60},
}

% Environnement exercice — argument optionnel = référence page du livre
\newcounter{exercice}[section]
\renewcommand{\theexercice}{\thesection.\arabic{exercice}}

\newtcolorbox[use counter=exercice]{exercice}[1][]{%
	colback=exo-bg,
	colframe=exo-border,
	fonttitle=\bfseries,
	title={Exercice~\theexercice \ifthenelse{\equal{#1}{}}{}{~---~#1}},
	breakable,
	before upper={\parindent=1em}
}

\newtcolorbox{solution}{%
	colback=sol-bg,
	colframe=sol-border,
	fonttitle=\bfseries,
	title={Solution},
	breakable,
	before upper={\parindent=1em}
}


%%%%%%%%%%%%%%%%%%%%%%%%%%%%%%%%%%%%%%%%%%%%%%%%%%%
%% Début du document
%%%%%%%%%%%%%%%%%%%%%%%%%%%%%%%%%%%%%%%%%%%%%%%%%%%
\begin{document}

\maketitle
\sloppy


%%%%%%%%%%%%%%%%%%%%%%%%%%%%%%%%%%%%%%%%%%%%%%%%%%%
\section{Exercices du Chapitre 4}
\label{sec:chap4}
%%%%%%%%%%%%%%%%%%%%%%%%%%%%%%%%%%%%%%%%%%%%%%%%%%%

% ----------------------------------------------------------------
\begin{exercice}[p.~130]
Prove that $\ket{\psi_0} = \bigotimes_{j=0}^{n-1} \ket{+}$ has the minimum possible energy for
\[
  H_0 = -\sum_{j=0}^{n-1} X_j,
\]
by first showing that, for each $j$ and each state $\ket{\psi}$, it holds that
$\langle \psi | X_j | \psi \rangle \leq 1$, and then showing that
$\langle \psi_0 | X_j | \psi_0 \rangle = 1$ for each $j$.
\end{exercice}

\begin{solution}
Soit $\ket{\psi}$ un état quantique normalisé à $n$ qubits. Nous pouvons le décomposer dans la base computationnelle sous la forme :
\[
\ket{\psi} = \sum_{i=0}^{2^n-1} a_i \ket{i}, \quad \text{avec} \quad \sum_{i=0}^{2^n-1} |a_i|^2 = 1.
\]
L'opérateur $X_j$ agit sur l'état de base $\ket{i}$ en inversant le $j$-ème bit de $i$. Algébriquement, cette opération de flip se traduit par la règle suivante :
\[
X_j \ket{i} = \ket{i + (-1)^{i_j}2^j}
\]
Puisque $X_j$ est une matrice de Pauli, elle est sa propre réciproque ($X_j^2 = I$). Ainsi, si l'on pose $\ket{k} = X_j\ket{i}$, alors on a de manière bijective $\ket{i} = X_j\ket{k}$, ce qui donne la relation réciproque sur l'indice : $i = k + (-1)^{k_j}2^j$.

Calculons maintenant la valeur moyenne $\langle \psi | X_j | \psi \rangle$ en développant la double somme :
\[
\langle \psi | X_j | \psi \rangle = \sum_{k=0}^{2^n-1} \sum_{i=0}^{2^n-1} a_k^* a_i \bra{k} X_j \ket{i} = \sum_{k=0}^{2^n-1} \sum_{i=0}^{2^n-1} a_k^* a_i \braket{k | i + (-1)^{i_j}2^j}
\]
Le produit scalaire de la base orthogonale $\braket{k | i + (-1)^{i_j}2^j}$ est non nul (égal à $1$) si et seulement si $k = i + (-1)^{i_j}2^j$, ou de manière équivalente $i = k + (-1)^{k_j}2^j$. En injectant cette condition, la double somme devient :
\[
\langle \psi | X_j | \psi \rangle = \sum_{k=0}^{2^n-1} a_k^* a_{k + (-1)^{k_j}2^j}
\]
Posons $b_k = a_{k + (-1)^{k_j}2^j}$. L'opération de flip n'étant qu'une permutation globale des indices, l'ensemble des coefficients $b_k$ est une réorganisation exacte des coefficients $a_k$. Par conséquent, la norme du vecteur des amplitudes reste inchangée : $\sum |b_k|^2 = \sum |a_k|^2 = 1$.

En appliquant l'inégalité de \textbf{Cauchy-Schwarz} sur le produit scalaire des vecteurs d'amplitudes $\vec{a}$ et $\vec{b}$, on obtient :
\[
\langle \psi | X_j | \psi \rangle = \sum_{k=0}^{2^n-1} a_k^* b_k \leq \sqrt{\sum_{k=0}^{2^n-1} |a_k|^2} \times \sqrt{\sum_{k=0}^{2^n-1} |b_k|^2}
\]
Comme l'état est normalisé, on remplace par les valeurs :
\[
\langle \psi | X_j | \psi \rangle \leq \sqrt{1} \times \sqrt{1} = 1
\]
On a donc bien $\langle \psi | X_j | \psi \rangle \leq 1$ pour tout état $\ket{\psi}$ et tout qubit $j$.

\subsection*{Étape 2 : Évaluation pour l'état fondamental $\ket{\psi_0}$}

Considérons à présent l'état $\ket{\psi_0} = \bigotimes_{j=0}^{n-1} \ket{+}$. Cet état correspond à la superposition uniforme de tous les états de la base computationnelle :
\[
\ket{\psi_0} = \frac{1}{\sqrt{2^n}} \sum_{k=0}^{2^n-1} \ket{k}
\]
Ici, pour tout $k$, l'amplitude est constante et vaut $a_k = \frac{1}{\sqrt{2^n}}$. En reprenant la formule simplifiée établie à l'étape précédente, on a :
\[
\langle \psi_0 | X_j | \psi_0 \rangle = \sum_{k=0}^{2^n-1} a_k a_{k + (-1)^{k_j}2^j} = \sum_{k=0}^{2^n-1} \left( \frac{1}{\sqrt{2^n}} \times \frac{1}{\sqrt{2^n}} \right) = \sum_{k=0}^{2^n-1} \frac{1}{2^n}
\]
La somme contenant exactement $2^n$ termes (les $2^n$ configurations de la base), on obtient :
\[
\langle \psi_0 | X_j | \psi_0 \rangle = 2^n \times \frac{1}{2^n} = 1
\]

\subsection*{Conclusion : Minimum de l'énergie}

D'après l'étape 1, l'énergie moyenne de n'importe quel état $\ket{\psi}$ pour l'Hamiltonien $H_0$ est minorée par :
\[
\langle \psi | H_0 | \psi \rangle = -\sum_{j=0}^{n-1} \langle \psi | X_j | \psi \rangle \geq -\sum_{j=0}^{n-1} 1 = -n
\]
L'énergie minimale théorique possible est donc $-n$. 

D'après l'étape 2, en évaluant l'énergie de l'état particulier $\ket{\psi_0}$, on trouve :
\[
\langle \psi_0 | H_0 | \psi_0 \rangle = -\sum_{j=0}^{n-1} \langle \psi_0 | X_j | \psi_0 \rangle = -\sum_{j=0}^{n-1} 1 = -n
\]
L'état $\ket{\psi_0}$ est donc bien l'état fondamental (\textit{ground state}) possédant le minimum possible d'énergie pour l'Hamiltonien $H_0$.
\end{solution}


% ----------------------------------------------------------------
\begin{exercice}[p.~135]
Create an instance of a simple QUBO problem and solve it with an annealer.
Notice that the values for the variables in the solution will be 0 and 1 instead of 1 and $-1$.
\end{exercice}

\begin{solution}
On choisit un problème quadratique binaire (QUBO) à trois variables $x_0, x_1, x_2 \in \{0, 1\}$ défini par l'expression d'énergie suivante :
\[
  E(x_0, x_1, x_2) = -x_1 + 2x_0 x_1 - 3x_0 x_2
\]
On le formule avec \texttt{dimod} en utilisant le type \texttt{dimod.BINARY} et on le résout par recuit simulé (\texttt{SimulatedAnnealingSampler}).

\begin{tcolorbox}[colback=gray!8, colframe=gray!40, title=Code Python, fonttitle=\bfseries\small, breakable]
\begin{lstlisting}[language=Python]
import dimod
from dwave.samplers import SimulatedAnnealingSampler

# Definition de la matrice QUBO (h pour le lineaire, J pour le quadratique)
J = {(0, 1): 2, (0, 2): -3}
h = {1: -1}

# On specifie le type BINARY pour travailler sur les valeurs {0, 1}
problem = dimod.BinaryQuadraticModel(h, J, 0.0, dimod.BINARY)

sampler = SimulatedAnnealingSampler()
result = sampler.sample(problem, num_reads=10)
print(result)
\end{lstlisting}
\end{tcolorbox}

L'exécution du script produit la table d'échantillonnage suivante :
\begin{verbatim}
   0  1  2 energy num_oc.
0  1  0  1    -3.0       1
1  1  0  1    -3.0       1
2  1  0  1    -3.0       1
3  1  0  1    -3.0       1
4  1  0  1    -3.0       1
5  1  0  1    -3.0       1
6  1  0  1    -3.0       1
7  1  0  1    -3.0       1
8  1  0  1    -3.0       1
9  1  0  1    -3.0       1
['BINARY', 10 rows, 10 samples, 3 variables]
\end{verbatim}

Les variables sont bien de type \texttt{BINARY} ($0/1$; la solution optimale est $(x_0, x_1, x_2) = (1, 0, 1)$ avec $E = -3.0$, état fondamental du QUBO.
\end{solution}


% ----------------------------------------------------------------
\begin{exercice}[p.~143]
Check that the QUBO problem returned by \texttt{cqm\_to\_bqm} coincides with what you
would obtain should you apply the transformations explained in Section~3.4.1.
Don't forget the offset!
\end{exercice}

\begin{solution}

On part du problème contraint (Section~3.4.1) :
\[
  \text{Minimiser}\quad -2y_0 - 3y_1 \qquad
  \text{s.c.}\quad y_0 + 2y_1 \leq 2,\quad y_j \in \{0,1\}.
\]

\textbf{Transformation manuelle.}
Le membre de gauche de la contrainte peut prendre des valeurs allant de $0$ à $3$ (si $y_0=1$ et $y_1=1$). Pour satisfaire l'inégalité $\leq 2$, l'écart maximal à combler est de $2$ (lorsque $y_0 = y_1 = 0$).

Ocean introduit donc deux variables de slack de poids unitaire, notées $s_0, s_1 \in \{0,1\}$, afin de pouvoir rajouter précisément $0, 1$ ou $2$ à l'équation. La contrainte d'égalité s'écrit alors :
\[
  y_0 + 2y_1 + s_0 + s_1 = 2
\]
En utilisant le multiplicateur de Lagrange $\lambda = 5$, le Hamiltonien pénalisé non contraint devient :
\[
  H = -2y_0 - 3y_1 + 5\bigl(y_0 + 2y_1 + s_0 + s_1 - 2\bigr)^2.
\]

En développant le carré à l'aide de l'identité $\left(\sum a_i\right)^2 = \sum a_i^2 + 2\sum_{i < j} a_i a_j$ et en exploitant la propriété des variables binaires ($v^2 = v$), on isole chaque terme :

\begin{center}
\renewcommand{\arraystretch}{1.3}
\begin{tabular}{lll}
\hline
Terme & Calcul algébrique & Valeur numérique \\
\hline
$y_0$ (linéaire) & $-2 + 5 \cdot \bigl(1^2 + 2 \cdot 1 \cdot (-2)\bigr) = -2 + 5(1 - 4)$ & $-17$ \\
$y_1$ (linéaire) & $-3 + 5 \cdot \bigl(2^2 + 2 \cdot 2 \cdot (-2)\bigr) = -3 + 5(4 - 8)$ & $-23$ \\
$s_0$ (linéaire) & $0 + 5 \cdot \bigl(1^2 + 2 \cdot 1 \cdot (-2)\bigr) = 5(1 - 4)$ & $-15$ \\
$s_1$ (linéaire) & $0 + 5 \cdot \bigl(1^2 + 2 \cdot 1 \cdot (-2)\bigr) = 5(1 - 4)$ & $-15$ \\
$y_0 y_1$ (quadratique) & $5 \cdot (2 \cdot 1 \cdot 2)$ & $20$ \\
$y_0 s_0$ (quadratique) & $5 \cdot (2 \cdot 1 \cdot 1)$ & $10$ \\
$y_0 s_1$ (quadratique) & $5 \cdot (2 \cdot 1 \cdot 1)$ & $10$ \\
$y_1 s_0$ (quadratique) & $5 \cdot (2 \cdot 2 \cdot 1)$ & $20$ \\
$y_1 s_1$ (quadratique) & $5 \cdot (2 \cdot 2 \cdot 1)$ & $20$ \\
$s_0 s_1$ (quadratique) & $5 \cdot (2 \cdot 1 \cdot 1)$ & $10$ \\
offset (constante) & $5 \cdot (-2)^2$ & $20$ \\
\hline
\end{tabular}
\end{center}

Ces coefficients calculés à la main coïncident avec la structure du modèle binaire renvoyée par \texttt{dimod.cqm\_to\_bqm} :

\begin{tcolorbox}[colback=gray!8, colframe=gray!40, title=Validation des Coefficients (Manuel vs Ocean SDK), fonttitle=\bfseries\small, breakable]
\begin{center}
\renewcommand{\arraystretch}{1.4}
\small
\begin{tabular}{llcc}
\hline
\textbf{Variable (Manuel)} & \textbf{Variable (Ocean SDK)} & \textbf{Valeur} & \textbf{Statut} \\
\hline
$y_0$ & \texttt{y0} & $-17.0$ & \checkmark \\
$y_1$ & \texttt{y1} & $-23.0$ & \checkmark \\
$s_0$ & \texttt{slack\_vafea44a25936459bad8a6f0de953a4ab\_0} & $-15.0$ & \checkmark \\
$s_1$ & \texttt{slack\_vafea44a25936459bad8a6f0de953a4ab\_1} & $-15.0$ & \checkmark \\
$y_0 y_1$ & \texttt{('y1', 'y0')} & $\phantom{-}20.0$ & \checkmark \\
$y_0 s_0$ & \texttt{('slack\_...\_0', 'y0')} & $\phantom{-}10.0$ & \checkmark \\
$y_0 s_1$ & \texttt{('slack\_...\_1', 'y0')} & $\phantom{-}10.0$ & \checkmark \\
$y_1 s_0$ & \texttt{('slack\_...\_0', 'y1')} & $\phantom{-}20.0$ & \checkmark \\
$y_1 s_1$ & \texttt{('slack\_...\_1', 'y1')} & $\phantom{-}20.0$ & \checkmark \\
$s_0 s_1$ & \texttt{('slack\_...\_1', 'slack\_...\_0')} & $\phantom{-}10.0$ & \checkmark \\
offset & \textit{Offset global} & $\phantom{-}20.0$ & \checkmark \\
\hline
\end{tabular}
\end{center}
\end{tcolorbox}

Les identifiants longs \texttt{slack\_vafea...\_0} et \texttt{slack\_vafea...\_1} correspondent à nos variables $s_0$ et $s_1$. L'offset global est bien de $20.0$, ce qui valide la correspondance entre transformation manuelle et tranformation automatique d'Ocean SDK.

\end{solution}


% ----------------------------------------------------------------
\begin{exercice}[p.~151]
Pick the qubits numbered from 0 to 7 and check that their connections correspond
to the description of the Chimera topology that we have made in the text. Notice
that they all lie in the top-left corner cell, so they will be connected to one cell on
the right and one cell below.
\end{exercice}

\begin{solution}

\subsection*{Topologie Chimera — Structure d'une cellule}

Dans la topologie Chimera, les qubits sont organisés en une grille de cellules, chaque cellule contenant 8 qubits répartis en deux groupes selon leur orientation physique :
\begin{itemize}
  \item \textbf{Groupe horizontal} : qubits $0, 1, 2, 3$
  \item \textbf{Groupe vertical} : qubits $4, 5, 6, 7$
\end{itemize}
Au sein d'une même cellule, chaque qubit horizontal croise les 4 qubits verticaux, formant un graphe biparti complet $K_{4,4}$. Il n'y a aucune connexion interne entre qubits de même orientation.

Pour les connexions inter-cellules (externes), la puce du \texttt{DW\_2000Q\_6} est configurée sous forme d'une grille de $16 \times 16$ cellules. Chaque cellule contenant 8 qubits, le saut d'indexation pour passer à la cellule en dessous est de $8 \times 16 = 128$, tandis que le saut pour passer à la cellule à droite est de $8$.

\subsection*{Vérification via \texttt{sampler.adjacency}}

L'analyse de l'attribut \texttt{sampler.adjacency} pour la première cellule (qubits 0 à 7) donne les voisinages exacts suivants :

\begin{center}
\renewcommand{\arraystretch}{1.3}
\begin{tabular}{cll}
\hline
\textbf{Qubit} & \textbf{Voisins détectés} & \textbf{Interprétation de la connexion externe} \\
\hline
0 & 4, 5, 6, 7, \textbf{128} & Vers l'homologue $0$ de la cellule du \textbf{bas} ($0 + 128$) \\
1 & 4, 5, 6, 7, \textbf{129} & Vers l'homologue $1$ de la cellule du \textbf{bas} ($1 + 128$) \\
2 & 4, 5, 6, 7, \textbf{130} & Vers l'homologue $2$ de la cellule du \textbf{bas} ($2 + 128$) \\
3 & 4, 5, 6, 7, \textbf{131} & Vers l'homologue $3$ de la cellule du \textbf{bas} ($3 + 128$) \\
4 & 0, 1, 2, 3, \textbf{12}  & Vers l'homologue $4$ de la cellule de \textbf{droite} ($4 + 8$) \\
5 & 0, 1, 2, 3, \textbf{13}  & Vers l'homologue $5$ de la cellule de \textbf{droite} ($5 + 8$) \\
6 & 0, 1, 2, 3, \textbf{14}  & Vers l'homologue $6$ de la cellule de \textbf{droite} ($6 + 8$) \\
7 & 0, 1, 2, 3, \textbf{15}  & Vers l'homologue $7$ de la cellule de \textbf{droite} ($7 + 8$) \\
\hline
\end{tabular}
\end{center}

Chaque qubit possède bien exactement 5 voisins : 4 liens internes ($K_{4,4}$) et 1 unique lien externe dirigé vers la cellule adjacente correspondante.

\subsection*{Représentation géométrique de la cellule}

Le schéma ci-dessous représente le croisement physique des qubits au sein de la cellule $(0,0)$ et leurs extensions vers les cellules voisines.

\begin{center}
\begin{tikzpicture}[
  hnode/.style={rounded corners=2pt, draw=blue!70, fill=blue!15, thick, minimum width=110pt, minimum height=10pt, font=\bfseries\small},
  vnode/.style={rounded corners=2pt, draw=red!70,  fill=red!15,  thick, minimum width=10pt, minimum height=110pt, font=\bfseries\small},
  ext/.style={draw=gray!60, dashed, thick, ->, >=stealth}
]
  % Qubits verticaux (4, 5, 6, 7)
  \node[vnode] (q4) at (1, 0) {}; \node[above, red!70] at (1, 2) {4};
  \node[vnode] (q5) at (2, 0) {}; \node[above, red!70] at (2, 2) {5};
  \node[vnode] (q6) at (3, 0) {}; \node[above, red!70] at (3, 2) {6};
  \node[vnode] (q7) at (4, 0) {}; \node[above, red!70] at (4, 2) {7};

  % Qubits horizontaux (0, 1, 2, 3)
  \node[hnode] (q0) at (2.5,  1) {}; \node[left, blue!70] at (0.4,  1) {0};
  \node[hnode] (q1) at (2.5,  0) {}; \node[left, blue!70] at (0.4,  0) {1};
  \node[hnode] (q2) at (2.5, -1) {}; \node[left, blue!70] at (0.4, -1) {2};
  \node[hnode] (q3) at (2.5, -2) {}; \node[left, blue!70] at (0.4, -2) {3};

  % Points d'intersection (coupleurs internes)
  \foreach \x in {1,2,3,4} {
    \foreach \y in {-2,-1,0,1} {
      \fill[gray!80] (\x,\y) circle (2pt);
    }
  }

  % Connexions externes vers le bas (0-3 -> 128-131)
  \draw[ext] (0.4,  1) -- (-0.8,  1) node[left, font=\scriptsize] {vers 128};
  \draw[ext] (0.4,  0) -- (-0.8,  0) node[left, font=\scriptsize] {vers 129};
  \draw[ext] (0.4, -1) -- (-0.8, -1) node[left, font=\scriptsize] {vers 130};
  \draw[ext] (0.4, -2) -- (-0.8, -2) node[left, font=\scriptsize] {vers 131};

  % Connexions externes vers la droite (4-7 -> 12-15)
  \draw[ext] (1, -2) -- (1, -3.2) node[below, font=\scriptsize] {vers 12};
  \draw[ext] (2, -2) -- (2, -3.2) node[below, font=\scriptsize] {vers 13};
  \draw[ext] (3, -2) -- (3, -3.2) node[below, font=\scriptsize] {vers 14};
  \draw[ext] (4, -2) -- (4, -3.2) node[below, font=\scriptsize] {vers 15};

\end{tikzpicture}
\end{center}

\end{solution}


%%%%%%%%%%%%%%%%%%%%%%%%%%%%%%%%%%%%%%%%%%%%%%%%%%%
%% Fin du document
%%%%%%%%%%%%%%%%%%%%%%%%%%%%%%%%%%%%%%%%%%%%%%%%%%%
\end{document}
